\documentclass[12pt,a4paper]{article}
\usepackage[margin=1in]{geometry}
\usepackage{amsmath}
\usepackage{amssymb}
\usepackage{hyperref}
\usepackage{fancyhdr}
\usepackage{listings}
\usepackage{xcolor}

\lstset{
    basicstyle=\ttfamily\footnotesize,
    keywordstyle=\color{blue},
    stringstyle=\color{red},
    commentstyle=\color{green!60!black},
    breaklines=true,
    frame=single,
    backgroundcolor=\color{gray!10}
}

\pagestyle{fancy}
\fancyhf{}
\rhead{Pendulum Simulation Documentation}
\lhead{Design Lab}
\cfoot{\thepage}

\title{\textbf{Interactive Multi-Link Pendulum Simulator}\\
\large Quick Reference Guide}
\author{Yash Kumar \\ 21CS30059}
\date{\today}

\begin{document}

\maketitle

\tableofcontents
\newpage

\section{Overview}

Interactive physics-based pendulum simulator (1-3 links). Uses eXtended Position-Based Dynamics (XPBD) for constraints, collision detection, and optional energy conservation.

\section{Core Physics}

Time-stepping with gravity and constraint solving. 50 substeps per frame at 1/60 s.

\subsection{Gravity and Position Update}

\begin{align}
\vec{v}_i &\leftarrow \vec{v}_i + (0, -10) \Delta t \\
\vec{x}_i &\leftarrow \vec{x}_i + \vec{v}_i \Delta t
\end{align}

\noindent\textbf{Implementation Snippet} (\texttt{pendulum.js}, line 340):
\begin{lstlisting}[language=JavaScript]
for(i=1;i<numPoints;i++){
    p=points[i];
    p.vel.y-=gravity*sdt;
    p.prev.assign(p.pos);
    p.pos.add(p.vel,sdt);
}
\end{lstlisting}

\subsection{Distance Constraints (XPBD)}

Links maintain fixed length $L$: \quad $\|\vec{x}_1 - \vec{x}_0\| = L$

Constraint force with Lagrange multiplier $\lambda$:
\begin{equation}
\lambda = -\frac{C(\vec{x})}{w_0 + w_1 + \alpha/\Delta t^2}
\end{equation}

Corrections (weighted by inverse mass $w=1/m$):
\begin{equation}
\vec{x}_i \leftarrow \vec{x}_i \pm w_i \lambda \vec{n}
\end{equation}

\noindent\textbf{Implementation Snippet} (\texttt{pendulum.js}, lines 310-322):
\begin{lstlisting}[language=JavaScript]
function solveDistPos(p0,p1,d0,compliance,unilateral,dt){
    var w=p0.invMass+p1.invMass;
    if(w==0) return;
    var grad=p1.pos.minus(p0.pos);
    var d=grad.normalize();
    w+=compliance/dt/dt;
    var lambda=(d-d0)/w;
    if(lambda<0&&unilateral) return;
    p1.force=lambda/dt/dt;
    p1.elongation=d-d0;
    p0.pos.add(grad,p0.invMass*lambda);
    p1.pos.add(grad,-p1.invMass*lambda);
}
\end{lstlisting}

\subsection{Damping}

Edge damping: reduces motion along links (0.0, 10.0, 100.0 Ns/m)

Global damping: reduces overall motion (0.0, 0.5, 1.0, 2.0 Ns/m)

\subsection{Collisions}

Elastic reflection: $\vec{v}_{\text{new}} = \vec{v} - 2(\vec{n} \cdot \vec{v})\vec{n}$

Detects point-obstacle and link-obstacle collisions.

\noindent\textbf{Implementation Snippet} (\texttt{pendulum.js}, velocities update on line 347):
\begin{lstlisting}[language=JavaScript]
var v=p.pos.minus(p.prev);
var vn=n.dot(v);
if(vn<0){
    v.add(n,-2*vn);
    p.prev.assign(p.pos.minus(v));
}
\end{lstlisting}

\subsection{Energy Conservation}

When enabled, rescales velocities to maintain:
\begin{equation}
E = \sum_i \left( m_i g y_i + \frac{1}{2}m_i v_i^2 \right) = \text{constant}
\end{equation}

\noindent\textbf{Implementation Snippet} (\texttt{pendulum.js}, lines 356-357):
\begin{lstlisting}[language=JavaScript]
function computeEnergy(){
    var E=0;
    for(i=1;i<numPoints;i++){
        p=points[i];
        E+=p.pos.y/p.invMass*gravity+0.5/p.invMass*p.vel.lenSquared();
    }
    return E;
}
function forceEnergyConservation(prevE){
    var dE=(computeEnergy()-prevE)/(numPoints-1);
    if(dE<0){
        var postE=computeEnergy();
        for(i=1;i<numPoints;i++){
            p=points[i];
            var Ek=0.5/p.invMass*p.vel.lenSquared();
            var s=Math.sqrt((Ek-dE)/Ek);
            p.vel.scale(s);
        }
    }
}
\end{lstlisting}

\section{Key Parameters}

\begin{center}
\small
\begin{tabular}{|l|l|}
\hline
\textbf{Parameter} & \textbf{Value/Range} \\
\hline
Substeps per frame & 50 \\
Frame time step & 1/60 s \\
Gravity & 10 m/s$^2$ \\
Default mass & 1.0 kg \\
Default link length & 0.3 m \\
Max links & 3 \\
Max obstacles & 10 \\
Edge damping & 0, 10, 100 Ns/m \\
Global damping & 0, 0.5, 1.0, 2.0 Ns/m \\
\hline
\end{tabular}
\end{center}

\section{Controls}

\subsection{Main Buttons}

\begin{itemize}
\item \textbf{Restart}: Reset to cascading position
\item \textbf{Equilibrium}: Reset to hanging straight down
\item \textbf{Pause/Step}: Toggle pause; step one frame when paused
\item \textbf{Run}: Resume simulation
\end{itemize}

\noindent\textbf{Implementation Snippet} (\texttt{pendulum.js}, lines 391-392):
\begin{lstlisting}[language=JavaScript]
function step() { paused = true; timeStep(); }
function run() { if (paused) { paused = false; timeStep(); } }
\end{lstlisting}

\subsection{Configuration}

\textbf{Number of Links} (1-3): Adjusts active links; updates mass table automatically.

\textbf{Mass Table}: For each link set mass (0.001-10 kg) and length (0.2-1.0 m). Checkbox enables unilateral (rope) constraint.

\noindent\textbf{Implementation Snippet} (\texttt{pendulum.js}, lines 41-46):
\begin{lstlisting}[language=JavaScript]
document.getElementById("segsSlider").oninput = function() {
    numPoints = Number(this.value) + 1;
    document.getElementById("numSegs").innerHTML = this.value;
    resetPos(false);
    updateMassTable();
}
\end{lstlisting}

\subsection{Physics Controls}

\begin{itemize}
\item \textbf{Edge Damping}: Reduces link oscillation
\item \textbf{Global Damping}: Causes system to settle
\item \textbf{Energy Conservation}: Maintains constant total energy
\item \textbf{Show Trail}: Red path of last point
\item \textbf{Show Forces}: Red force, blue acceleration vectors
\end{itemize}

\subsection{Obstacles}

Add 1-10 circular obstacles. Click and drag to reposition. Turn orange when dragged.

\subsection{Mouse Interaction}

Click and drag pendulum points or obstacles to control them manually.

\noindent\textbf{Implementation Snippet} (\texttt{pendulum.js}, lines 464-467):
\begin{lstlisting}[language=JavaScript]
canvas.addEventListener("mousedown", onMouse);
canvas.addEventListener("mousemove", onMouse);
canvas.addEventListener("mouseup", onMouse);
canvas.addEventListener("mouseout", onMouse);
\end{lstlisting}

\section{Quick Examples}

\subsection{Simple Pendulum}

Links: 1, disable damping, enable energy conservation. Click Restart.

\subsection{Double Pendulum Chaos}

Links: 2 (0.4 m, 0.3 m), disable damping. Run multiple times with small position variations.

\subsection{Damped Motion}

Links: 2, global damping 2.0. Compare settling with vs. without damping.

\section{Technical Details}

Per substep: apply gravity → satisfy constraints → detect collisions → apply damping → recompute velocities. 50 substeps per frame.

Canvas: 500$\times$500 px representing 2.0 m wide world. Scaling: 250 px/meter.

\section{Troubleshooting}

\begin{itemize}
\item \textbf{Too slow}: Reduce obstacles
\item \textbf{Points escape}: Check mass values; reduce compliance
\item \textbf{Constraints elastic}: Reduce damping
\item \textbf{Bad collisions}: Increase damping or enable energy conservation
\end{itemize}

\section{Conclusion}

Experiment with different configurations to explore classical mechanics, chaos, damping, constraints, and energy conservation.

\end{document}